% =============================================================================
\chapter{Methodology}
\label{ch:methodology}
% =============================================================================

This chapter explains how the data described in
Chapter~\ref{ch:data} are used in the regressions that test the
two hypotheses set out in Chapter~\ref{ch:literature}.  Section
\ref{sec:design} gives an overview of the empirical strategy.
Sections~\ref{sec:variables} and \ref{sec:models} define the regression
variables and construct three nested regression equations.
Section~\ref{sec:estimation} describes how the models are estimated
and how their standard errors are computed.  Section~\ref{sec:hypothesis_tests} states the formal decision rules
used to evaluate the results for H1 and H2.  Finally, section~\ref{sec:robustness_plan}
describes the robustness checks reported in
Chapter~\ref{ch:robustness}.


% ─────────────────────────────────────────────────────────────────────────────
\section{Empirical strategy}
\label{sec:design}
% ─────────────────────────────────────────────────────────────────────────────

The thesis estimates three panel regressions of post-filing
stock-return volatility on increasingly rich sets of explanatory
variables.  The \emph{baseline} model uses lagged volatility together
with standard accounting controls.  The \emph{text} model adds three
features extracted from the 10-K narrative.  The \emph{divergence}
model further adds a single composite measure that captures the gap
between the financial and the narrative trajectory of a firm from one
year to the next.

Hypothesis~H1 holds that the textual features explain volatility
beyond what financial fundamentals already do.  It is supported if the
text model improves on the baseline both in adjusted~$R^2$ and in a
joint $F$-test on the textual coefficients.  Hypothesis~H2 holds that
the divergence measure carries information beyond \emph{both}
financials and the textual features taken individually.  It is
supported if the divergence coefficient is positive and statistically
significant, and if the divergence model improves on the text model.

The three-step nesting is deliberate.  By first absorbing the
predictable, mechanical part of volatility (persistence and firm
fundamentals) and only then asking whether language adds further
explanatory power, the design avoids attributing to text variation
that any reasonable accounting model could already account for.


% ─────────────────────────────────────────────────────────────────────────────
\section{Variables}
\label{sec:variables}
% ─────────────────────────────────────────────────────────────────────────────

\subsection{Dependent variable: post-filing volatility}
\label{ssec:dependent_variable}

The dependent variable is the \emph{annualised standard deviation of
daily log returns} computed over a 365-calendar-day window that begins
two trading days after the 10-K filing date.  The two-day lag avoids
any same-day price reaction to the filing itself, while the one-year
window matches the annual frequency of the panel and reflects the
horizon over which a 10-K is typically the most recent piece of
public information.  Formally,
\begin{equation}
\sigma_{i,t+1} \;=\; \sqrt{252}\;\cdot\;
\operatorname{sd}\!\bigl(r_{i,d}\bigr)_{d \in W_{i,t}},
\label{eq:vol}
\end{equation}
where $r_{i,d} = \ln(P_{i,d}/P_{i,d-1})$ is the daily log return for
firm~$i$ and $W_{i,t}$ is the post-filing window for filing year~$t$.
The factor $\sqrt{252}$ rescales the daily standard deviation to an
annual figure, which is the conventional unit in the volatility
literature.  A firm-year is kept only if the window contains at least
200~trading days, and a firm enters the panel only if at least 90\%
of its expected windows meet this requirement.

\subsection{Financial controls}
\label{ssec:financial_controls}

The baseline regressors are taken directly from the firm's annual
XBRL filings (Section~\ref{sec:xbrl_tags}).  Five controls enter the
full baseline:
\emph{lagged volatility} (the same measure as the dependent variable
but for the previous filing year), \emph{firm size} measured as the
natural logarithm of total assets, \emph{leverage} (total liabilities
over total assets), \emph{return on assets} (net income over total
assets), and \emph{asset growth} (the year-on-year change in total
assets divided by lagged total assets).  All ratios are winsorised at
the 1st and 99th percentiles within the regression sample to reduce
the influence of extreme observations and small-denominator artefacts.

\subsection{Textual features}
\label{ssec:textual_features}

Three textual features summarise the narrative content of each 10-K:
\begin{itemize}
  \item $\Delta\text{Sentiment}_{i,t}$: the year-on-year change in
        net tone of Item~7 (MD\&A), measured with the
        \textcite{loughran2011} positive and negative word lists.
        Net tone is the count of positive minus negative words,
        scaled by the total word count of the section.
  \item $\Delta\text{Risk}_{i,t}$: the year-on-year change in the
        share of risk-related words in Item~1A (Risk Factors), using
        the same dictionary's uncertainty and litigious lists.
  \item $\text{TextSim}_{i,t}$: the cosine similarity between the
        full 10-K text in year~$t$ and year~$t-1$, computed on
        TF--IDF vectors built from the cleaned narrative
        (Section~\ref{ssec:text_postprocessing}).  Higher values
        mean the firm's disclosure is more similar to the previous
        year; lower values indicate substantive changes in language.
\end{itemize}
All three features are by construction first differences (or
similarity to the previous filing) and therefore capture changes in
narrative content rather than cross-sectional levels, which would
mostly reflect industry vocabulary.

\subsection{Divergence}
\label{ssec:divergence}

The central novel variable is \emph{Divergence}, a one-dimensional
measure of the gap between the financial and the narrative
trajectory of a firm.  It is constructed as the firm-demeaned
residual of $\Delta\text{Sentiment}$ on $\Delta\text{ROA}$:
\begin{equation}
\text{Divergence}_{i,t} \;=\;
\bigl(\Delta\text{Sentiment}_{i,t} - \overline{\Delta\text{Sentiment}}_{i}\bigr)
\;-\; \lambda \cdot
\bigl(\Delta\text{ROA}_{i,t} - \overline{\Delta\text{ROA}}_{i}\bigr),
\label{eq:divergence}
\end{equation}
where $\Delta\text{ROA}_{i,t} = \text{ROA}_{i,t} - \text{ROA}_{i,t-1}$,
$\overline{X}_{i}$ denotes the within-firm mean of $X$, and $\lambda$
is the slope coefficient of a within-firm regression of
$\Delta\text{Sentiment}$ on $\Delta\text{ROA}$ estimated on the
training sample.  Centring on the firm mean removes the firm-specific
intercept, so the construction is exactly the residual of the
within-firm regression and is therefore orthogonal to
$\Delta\text{ROA}$ within firm by construction.  The variable inherits
the units of $\Delta\text{Sentiment}$ and has a simple interpretation:
$\lambda \cdot \Delta\text{ROA}_{i,t}$ is the change in tone that
would, on average, accompany a change in profitability of size
$\Delta\text{ROA}_{i,t}$ in a given firm; subtracting it from the
actual change in tone gives the part of the narrative shift that is
\emph{not} explained by contemporaneous financial performance.  Large
positive values flag firms whose tone improved more than their
fundamentals warranted; large negative values flag the opposite.  As
a robustness check we also use a binary dummy that equals one when
$\operatorname{sign}(\Delta\text{Sentiment})$ differs from
$\operatorname{sign}(\Delta\text{ROA})$.


% ─────────────────────────────────────────────────────────────────────────────
\section{Model specifications}
\label{sec:models}
% ─────────────────────────────────────────────────────────────────────────────

The three nested specifications are estimated on the same firm-year
panel with the same fixed-effect structure.  Let
$\mathbf{X}^{\text{fin}}_{i,t}$ collect the four financial controls
(size, leverage, ROA, asset growth) and let
$\mathbf{X}^{\text{text}}_{i,t}$ collect the three textual features.

\paragraph{Baseline model.}
\begin{equation}
\sigma_{i,t+1} \;=\; \alpha + \phi \, \sigma_{i,t}
   \;+\; \boldsymbol{\beta}' \mathbf{X}^{\text{fin}}_{i,t}
   \;+\; \mu_j \;+\; \tau_t \;+\; \varepsilon_{i,t},
\label{eq:baseline}
\end{equation}
where $\mu_j$ is a two-digit SIC industry fixed effect and $\tau_t$
is a fiscal-year fixed effect.  The model isolates the share of
volatility that is mechanically traceable to past volatility and
balance-sheet fundamentals.

\paragraph{Text model.}
\begin{equation}
\sigma_{i,t+1} \;=\; \alpha + \phi \, \sigma_{i,t}
   \;+\; \boldsymbol{\beta}' \mathbf{X}^{\text{fin}}_{i,t}
   \;+\; \boldsymbol{\gamma}' \mathbf{X}^{\text{text}}_{i,t}
   \;+\; \mu_j \;+\; \tau_t \;+\; \varepsilon_{i,t}.
\label{eq:text}
\end{equation}
Hypothesis~H1 corresponds to a joint test that
$\boldsymbol{\gamma} \neq \mathbf{0}$.

\paragraph{Divergence model.}
\begin{equation}
\sigma_{i,t+1} \;=\; \alpha + \phi \, \sigma_{i,t}
   \;+\; \boldsymbol{\beta}' \mathbf{X}^{\text{fin}}_{i,t}
   \;+\; \boldsymbol{\gamma}' \mathbf{X}^{\text{text}}_{i,t}
   \;+\; \delta \cdot \text{Divergence}_{i,t}
   \;+\; \mu_j \;+\; \tau_t \;+\; \varepsilon_{i,t}.
\label{eq:divergence_model}
\end{equation}
Hypothesis~H2 corresponds to $\delta > 0$ together with a positive
incremental adjusted~$R^2$ over the text model.


% ─────────────────────────────────────────────────────────────────────────────
\section{Estimation and inference}
\label{sec:estimation}
% ─────────────────────────────────────────────────────────────────────────────

All three models are estimated by ordinary least squares with
two-digit SIC industry and fiscal-year fixed effects absorbed using
the absorbing-regression routine of the
\texttt{linearmodels} package.  Coefficients are therefore identified
from \emph{within-industry, within-year} variation, which removes the
risk that the results are driven by either time trends in
volatility (e.g.\ the 2008 and 2020 crises) or by persistent
industry-level differences in risk.

\paragraph{Standard errors.}
Standard errors are clustered at the firm level following
\textcite{petersen2009}, which makes inference robust to arbitrary
within-firm correlation.  Two features of the data make this
particularly important.  First, the dependent variable is itself a
volatility estimate computed on a 365-day window, so consecutive
firm-year windows from a single firm overlap by roughly two-thirds of
their length.  This induces strong serial correlation in
$\sigma_{i,t+1}$ that ordinary heteroskedasticity-robust standard
errors would understate.  Second, omitted firm-specific risk
characteristics are likely to persist over time.  Firm clustering
addresses both concerns simultaneously without imposing any
particular correlation structure.

\paragraph{Outliers.}
The financial ratios used as regressors are winsorised at the 1st and
99th percentiles of the regression sample
(Section~\ref{ssec:financial_controls}).  Winsorisation is preferred
to outright deletion because many of the extreme values are
economically real (e.g.\ acquisition-driven asset growth) but would
otherwise dominate OLS estimation through the squared-residual
objective.

\paragraph{Sample composition.}
After applying all data filters described in Chapter~\ref{ch:data}
and dropping firm-years with missing values for any regressor, the
estimation sample consists of about 6{,}300 firm-years on 537~unique
firms.  Twelve firms changed their fiscal year-end during the sample
period; the transition year is dropped, leaving a multi-year gap in
the time series of those firms.  Asset growth is set to missing at
the gap, while lagged volatility and lagged returns are unaffected
because they are anchored to filing dates rather than fiscal periods.


% ─────────────────────────────────────────────────────────────────────────────
\section{Testing the hypotheses}
\label{sec:hypothesis_tests}
% ─────────────────────────────────────────────────────────────────────────────

The two hypotheses are evaluated through a combination of incremental
fit statistics and formal tests on the relevant coefficients.

\paragraph{H1: text adds information beyond financials.}
The text model
(equation~\ref{eq:text}) is compared to the baseline
(equation~\ref{eq:baseline}) on two criteria:
\begin{enumerate}
  \item the text model's adjusted~$R^2$ exceeds the baseline's, and
  \item a Wald-style joint test on
        $H_0\!: \boldsymbol{\gamma} = \mathbf{0}$ rejects at the
        conventional 5\% level using firm-clustered standard errors.
\end{enumerate}
Both conditions must be met for H1 to be considered supported.

\paragraph{H2: divergence adds information beyond text.}
The divergence model (equation~\ref{eq:divergence_model}) is
compared to the text model on two criteria:
\begin{enumerate}
  \item the divergence coefficient $\delta$ is positive and
        significant at the 5\% level, and
  \item the divergence model's adjusted~$R^2$ exceeds the text
        model's.
\end{enumerate}
A positive sign is required because the theoretical prediction is
directional: if narrative tone has drifted away from concurrent
financial performance, future volatility should rise, not fall.

\paragraph{Economic magnitude.}
Statistical significance alone is not enough to justify a finding,
so the divergence coefficient is also reported in standardised form:
the change in $\sigma_{i,t+1}$ implied by a one-standard-deviation
increase in Divergence, expressed as a percentage of the sample
mean of the dependent variable.  In addition, the divergence sample
is split into quartiles and the average post-filing volatility of
each quartile is reported, which gives a model-free check that the
relationship is monotone rather than driven by the tails.


% ─────────────────────────────────────────────────────────────────────────────
\section{Planned robustness checks}
\label{sec:robustness_plan}
% ─────────────────────────────────────────────────────────────────────────────

Chapter~\ref{ch:robustness} reports a battery of checks intended to
make sure the main results do not depend on a single modelling
choice.  Briefly, we re-estimate the divergence model
\begin{enumerate*}[label=(\roman*)]
  \item with the binary divergence dummy in place of the continuous
        measure;
  \item with post-filing return as the dependent variable instead of
        volatility;
  \item excluding the 2020 fiscal year, which is heavily influenced
        by the COVID-19 shock;
  \item excluding firms whose 10-K extraction was flagged in the
        AI-assisted audit (Section~\ref{ssec:text_coverage}), to
        confirm that the result is not driven by parsing artefacts;
  \item with firm fixed effects in place of industry fixed effects,
        which is a stricter specification but reduces the effective
        sample.
\end{enumerate*}
Each check is reported in full in the tables of
Chapter~\ref{ch:robustness}.
